% Transcription — fonds Alexandre Grothendieck, shelfmark 23, batch 5 (pages 81-98).
% Established from the facsimile archives/batches/23/batch-05.pdf (a mirror of
% https://grothendieck.umontpellier.fr/23.pdf), per the skill
% transcribe-grothendieck. The batch file carries no cover sheet and opens at
% page 81: PDF page N is archivists' page 80+N. Checked against the pencilled
% numbers bottom-left — 81 on the orange cover, 82, 83 … 98 on the leaves that
% follow, 89 and 92-98 read clearly, the others faint but present — the
% alignment holds end to end. This is the folder's last batch: eighteen leaves,
% not twenty.
%
% Pass: Fable 5.1 (claude-fable-5-1), 2026-09-21 — first pass, unchecked against the pages by a human.
%
% What the batch is. Two orange cover sheets, each inscribed in his hand at the
% top right with a sub-dossier number of the chapter, cut the eighteen leaves
% into two runs of loose notes:
%
%   81      Orange cover, « IV 5 », then a struck word and « Compléments
%           d'ordre divers » in pencil. Its words open the first section; no
%           \page.
%   82-83   Ink, one leaf recto-verso: « Rappels sur la dimension » — a lemma
%           on the codimension of Y in V(x) for a section x of an invertible
%           Module, with two corollaries on systems of parameters, and a
%           further corollary written up the left margin of 83.
%   84      Ink: « ∂-Catégorie » — the first axiom of a graded additive
%           category with translation, X(1) and X(-1). The leaf is cancelled by
%           diagonal strokes; transcribed, the cancellation recorded once.
%   85      Blank (bleed-through of 84 only). Absent.
%   86-88   One argument written out of page order. The « Remarque » on the
%           finiteness of the integral closure of a quotient of a regular ring
%           begins on 88 and continues on 87; the margin of 87 says « cf verso »
%           and the verso is 86, which carries the lemma on the (S_k) property
%           of a finite algebra over A it relies on. The bottom of 87 is the
%           first, hatched-out draft of that lemma. Pages are transcribed in
%           archivists' order with a \note giving the reading order.
%   89      Ink draft over a pencil underlayer. The ink is scratch work for the
%           lemma of 86 (prof ≥ Inf(k, dim), the module against the ring);
%           the pencil, in English and partly under the ink, is a pair of
%           Lefschetz-type conditions on Pic(X') → Pic(Y') in terms of depth.
%           Both are transcribed, the layer said once.
%   90      Orange cover, « IV 12 », then in pencil « Fibres des morphismes
%           plats de présentation finie ». Its words open the second section;
%           no \page.
%   91-93   Ink over pencil, a fast draft: « Un exemple » — a local scheme X
%           glued from X' and X'' along a curve D, with a function f such that
%           the non-Cohen-Macaulay locus of V(f) has the wrong codimension;
%           then (93) the recipe for the construction: the cone over a
%           projectively normal non-Cohen-Macaulay S', an affine space, and
%           a spectrum of k[t] localised.
%   94      Blank (bleed-through of 93). Absent.
%   95-98   Blue ink, fair copy, one stapled run: a Proposition on graded
%           algebras A → B → C over K with a K-linear section of φ, its
%           conditions b), b'), b''), the Poincaré series P_{B,x} ≪ P_{A,x} P_{C,x},
%           two corollaries and a « Question » on flatness.
%
% Notation, fixed once for the batch. His « prof » is kept as \operatorname{prof},
% his « Codim », « ht », « Inf », « Sup », « rg » likewise; « C.M. », « Coh.
% Mac. » and « Cohen-Macaulay » are kept as he writes each of them. Prime ideals
% are his script p, q, m, set as \mathfrak. On pages 86-87 the letters p and q
% change sides between the statement of the lemma (q, q' primes of B over p of
% A) and its proof (p a prime of B over q of A); the swap is on the leaf and is
% transcribed as written, flagged once. The pencil drawings on page 86 (a
% child's hand, not his) are not transcribed.
\documentclass[11pt,a4paper]{article}
% Le préambule commun à toutes les transcriptions du fonds Grothendieck.
%
% Il tient en une idée : **l'incertitude se compose, elle ne se résout pas.**
% Une transcription qui lisse un mot douteux a détruit la seule chose pour
% laquelle on la faisait — un lecteur doit pouvoir savoir, à chaque endroit,
% ce qui était sur le feuillet et ce qui a été ajouté. Les six macros qui
% suivent sont donc l'appareil critique, et non de la décoration.
%
% Les mêmes commandes sont reconnues par `scripts/render.mjs`, qui produit la
% vue de lecture du site. Une macro ajoutée ici doit l'être aussi là-bas,
% faute de quoi le rendu échouera bruyamment — ce qui est voulu : un rendu
% silencieusement faux est pire qu'un rendu absent.

\ProvidesPackage{grothendieck}[2026/08/08 Transcriptions du fonds Grothendieck]

\usepackage{amsmath,amssymb,amsthm}
\usepackage{mathrsfs}
\usepackage{stmaryrd}
% Les diagrammes commutatifs sont partout dans ce fonds : c'est la langue
% même de la géométrie algébrique de Grothendieck, et une transcription qui
% les rendrait en prose ne transcrirait rien.
\usepackage{tikz-cd}
% Français par défaut — c'est la langue des feuillets — mais l'anglais est
% chargé aussi : la traduction et le résumé sont des documents anglais, et la
% typographie française (espaces avant les deux-points) y serait une faute.
% Ces documents basculent par \selectlanguage{english} en tête de corps.
\usepackage[english,french]{babel}
\usepackage{fontspec}
\usepackage{geometry}
\geometry{a4paper,margin=25mm,marginparwidth=28mm}
\usepackage{marginnote}
\usepackage{xcolor}
% ulem, not soul. `soul`'s \st and \ul are built on a character-by-character
% scan that XeLaTeX's font machinery breaks: the document fails with an
% undefined control sequence rather than degrading. ulem does the same two jobs
% and is Unicode-engine safe. `normalem` stops it hijacking \emph, which the
% transcriptions use for Grothendieck's own emphasis.
\usepackage[normalem]{ulem}
\usepackage{hyperref}
\hypersetup{colorlinks=true,linkcolor=black,urlcolor=black,pdfborder={0 0 0}}

% --- Métadonnées du lot -----------------------------------------------------
% Reprises telles quelles par le script de rendu, qui les affiche en tête de
% la vue de lecture. Elles fixent ce que le fichier prétend couvrir : sans
% elles, une transcription flotte sans point d'attache dans un fonds de seize
% mille feuillets.

\newcommand{\folder}[1]{\gdef\grothFolder{#1}}
\newcommand{\batch}[1]{\gdef\grothBatch{#1}}
\newcommand{\leaves}[2]{\gdef\grothFirst{#1}\gdef\grothLast{#2}}
\newcommand{\foldertitle}[1]{\gdef\grothFoldertitle{#1}}
\gdef\grothFolder{?}\gdef\grothBatch{?}\gdef\grothFirst{?}\gdef\grothLast{?}\gdef\grothFoldertitle{}

% --- Le résumé --------------------------------------------------------------
%
% Il ouvre la lecture modernisée, sous le titre « Résumé », et il est composé
% en petit. Ce n'est pas un ornement : c'est le seul passage du document qui
% ne soit pas des mathématiques, et un lecteur doit pouvoir le reconnaître
% avant de l'avoir lu — puis le sauter s'il connaît déjà le sujet, ou s'y
% arrêter s'il ne le connaît pas. Le filet qui le suit marque la reprise du
% corps.

\newenvironment{resume}
  {\small\setlength{\parskip}{0.45em}}
  {\par\medskip\hrule\medskip}

% --- Mots-clés ---------------------------------------------------------------
%
% En anglais, en clôture du résumé de la lecture modernisée : le vocabulaire
% moderne sous lequel chercher ce que les feuillets construisent. Le manifeste
% du site (`npm run manifest`) les extrait pour étiqueter la cote entière —
% cette ligne est la source unique des tags, il n'y a pas de fichier à côté.
\newcommand{\keywords}[1]{\par\smallskip\noindent
  {\footnotesize\textsc{Keywords} --- \textit{#1}}\par}

% --- L'appareil critique ----------------------------------------------------

% Le numéro de feuillet, en marge. C'est l'ancre qui permet de régler un
% désaccord sur une formule en regardant *un* feuillet plutôt qu'une chemise
% de six cents. Le site s'en sert aussi pour tourner les pages du fac-similé
% quand on fait défiler la transcription.
\newcommand{\leaf}[1]{%
  \marginnote{\footnotesize\color{gray!70}\textsf{#1}}%
  \label{leaf:#1}\ignorespaces}

% Illisible : on ne devine pas. Un mot inventé qui se lit comme les autres est
% le pire résultat possible.
\newcommand{\ill}{\textcolor{red!60!black}{[\,\ldots\,]}}

% Lecture douteuse : le mot est donné, et signalé comme incertain.
\newcommand{\uncertain}[1]{\uline{#1}}

% Ajout de l'éditeur — une lettre manquante, un mot sous-entendu.
\newcommand{\add}[1]{\textcolor{blue!50!black}{[#1]}}

% Biffé par Grothendieck lui-même. Ce qu'il a rayé fait partie du document :
% une rature dit souvent où la pensée a changé de direction.
\newcommand{\struck}[1]{\sout{#1}}

% Note du transcripteur, sur l'état matériel du feuillet.
\newcommand{\note}[1]{\par\smallskip{\small\color{gray!60!black}\textit{[#1]}}\par\smallskip}

% Note marginale de Grothendieck, distincte du corps du texte.
\newcommand{\marginal}[1]{\par\smallskip\noindent\hspace{1em}%
  {\small\color{gray!60!black}#1}\par\smallskip}

% --- Titre ------------------------------------------------------------------

\AtBeginDocument{%
  \begin{center}
    {\large\bfseries Fonds Alexandre Grothendieck --- cote n\textsuperscript{o}~\grothFolder}\\[2pt]
    {\itshape\grothFoldertitle}\\[4pt]
    {\small feuillets \grothFirst--\grothLast\ (lot \grothBatch)}\\[2pt]
    {\footnotesize\color{gray!60!black}Transcription établie d'après le fac-similé numérisé
     par l'Université de Montpellier.\\
     Les lectures incertaines sont soulignées, les illisibles notées [\,\ldots\,],
     les ajouts entre crochets.}
  \end{center}
  \vspace{1em}}

\endinput


\folder{23}
\batch{5}
\pages{81}{98}
\dating{1966-1968}
\watermark{Édition de démonstration}
\foldertitle{EGA IV. Compléments : notes manuscrites (s.d.), lettres (1966, 1968), tapuscrit (s.d.).}

\begin{document}

\section{IV 5. Compléments d'ordre divers}
\note{les mots sont de sa main, en haut à droite de la chemise orange qui
forme le feuillet 81 : « IV 5 », puis un mot biffé qui ne se lit plus, puis
« Compléments d'ordre divers ». La chemise ne porte rien d'autre.}

\page{82}

\marginal{N.B. \uncertain{dans} \struck{\ill{}} Rappels sur la dimension}
\note{encadré, en haut à droite du feuillet.}

\textbf{Lemme 1.} \struck{$A$ anneau local} $X$ préschéma \struck{\ill{}}
nœth., $Y$ partie fermée $\neq \emptyset$, $(Y_i)_{i \in I}$ ens. des comp.
irréductibles de $Y$, $(X_k)_{k \in K}$ ens. des comp. irréductibles de $X$.
\struck{$\exists$} $r = \operatorname{Codim}(Y, X)$, $I'$ ens. des $i \in I$
tels que $\operatorname{Codim}(Y_i, X) = r$. Pour $i \in I'$, $K_i$ = ens. des
$k \in K$ tels que \struck{Codim} $X_k \supset Y_i$ et
$\operatorname{Codim}(Y_i, X_k) = r$. Soit $x$ une section d'un Module inv.
$L$ sur $X$ telle que $V(x) \supset Y$. Alors :

1°) \struck{On a} \struck{\ill{}}
$\operatorname{Codim}(Y, V(x)) \geq r - 1$, et pour qu'on ait égalité, il f. et
s. qu'il existe $i \in I'$, $k \in K_i$ tels que $V(x) \not\supset X_k$.
\note{le signe d'inégalité est biffé et réécrit au-dessus ; en interligne
au-dessus de « tels que » : « $\operatorname{codim}(V(x), X) \geq 1$ i.e.
$= 1$ \ill{} ».}

2°) \struck{Supposons $r > 0$. Pour \ill{} \ill{} \ill{}} Supposons que $X$
soit \uncertain{local}, \ill{} \struck{\ill{} $A_x$, $X = V(\ill{})$} et
$r > 0$, alors il existe un \struck{\ill{}} $x \in \mathfrak{m}$ tel que
$\operatorname{Codim}(Y, V(x)) = r - 1$ \struck{$\dim V(x) = \dim X - 1$}
\ill{} \ill{}. $X$ est local, on \ill{} \ill{} $\dim V(x) = \dim X - 1$.
\note{la conclusion de 2°) est écrite en accolade, sur deux lignes dont la
seconde est biffée puis reprise ; la fin de la phrase ne se lit pas.}

Dém. 1) $\operatorname{Codim}(Y, V(x)) = \operatorname{Inf}_i \struck{\ill{}}
\operatorname{Codim}(Y_i, V(x))$, or
\[
  \operatorname{Codim}(Y_i, V(x)) = \dim \mathcal{O}_{V(x), y_i}
  \geq \dim \mathcal{O}_{X, y_i} - 1 = \operatorname{Codim}(Y_i, X) - 1
  \geq \operatorname{Codim}(Y, X) - 1 ,
\]
d'où $\operatorname{Codim}(Y, V(x)) \geq \operatorname{Codim}(Y, X) - 1$,
égalité ssi $\exists i$ tel qu'on \struck{\ill{}} ait les deux égalités
\[
  \begin{cases}
    \operatorname{Codim}(Y_i, V(x)) = \struck{\dim} \operatorname{Codim}(Y_i, X) - 1 \\
    \operatorname{Codim}(Y_i, X) = \operatorname{Codim}(Y, X)
  \end{cases}
\]
\note{« $\exists i$ tel qu'on » est entouré sur le feuillet.}

\page{83}

i.e. $i \in I'$, et $\exists k \in K_i$ t.q. $V(x) \not\supset X_k$.

2°) Comme $r > 0$, pour toutes les composantes irréd. $X_\alpha$ de $X$, on a
$X_\alpha \not\subset Y$, donc \struck{il existe $x \in \mathfrak{m}$ \ill{}}
\struck{Comme \ill{} \ill{} \ill{}} $\mathfrak{m} \not\subset
\mathfrak{p}_\alpha$, donc $\exists x \in \mathfrak{m}$, \struck{$I' \neq
\emptyset$ et} $x \notin \bigcup_\alpha \mathfrak{p}_\alpha$.

\struck{pour $i \in I'$, $K_i \neq \emptyset$, pour \ill{} \ill{} $r > 0$ ;
$X_k \supset Y_i$, $X_k \neq Y_i$, donc $X_k \not\subset Y$ ; donc \ill{}
$\mathfrak{p}_k \not\supset \mathfrak{m}$, donc $\exists x \in \mathfrak{m}$
t.q. $V(x)$ \ill{} \ill{} $\exists x \in \mathfrak{m}$, $x \notin
\mathfrak{p}_k$. \ill{} $\mathfrak{m} \supset \mathfrak{p}$, donc $\exists x
\in \mathfrak{m}$, $x \notin \mathfrak{p}$. \uncertain{Et} $\exists
\mathfrak{p}_\alpha$, pour toutes composantes irréductibles $X_\alpha$, \ill{}
on a $X_\alpha \not\subset Y$ [sinon}
\note{tout ce bloc, une première rédaction de 2°), est barré de traits
diagonaux ; il est remplacé par les trois lignes qui le précèdent. En marge
de gauche, à sa hauteur : « Je \uncertain{modifie} \ill{} conditions
\uncertain{indiquées} \ill{} \ill{} ».}

\textbf{Corollaire 1 :} Soit $A$ un anneau local nœth., \struck{$X =
\operatorname{Spec} A$, $Y = V(\mathfrak{m})$ \ill{} \ill{} $\mathfrak{m}$
\ill{} idéal, $\mathfrak{p}$ partie fermée de $X$} \struck{$x \in
\mathfrak{m}$} $r = \struck{\operatorname{Codim}(Y, X)}
\operatorname{ht}(\mathfrak{m})$, alors il existe une suite $x_1, \dots, x_r$
d'éléments de $\mathfrak{m}$, \struck{faisant partie d'un système de
paramètres de $A$, et tel que $\operatorname{ht}(x_1, \dots, x_r)$}
$\operatorname{ht}(\sum x_i A) = r$, d'où \struck{et}
$\operatorname{ht}(\mathfrak{m} / (x_1, \dots, x_r)) = 0$ ; si $A$ local, les
$x_i$ font partie d'un syst. de paramètres.
\note{« d'où » est entouré. L'énoncé hésite entre un anneau local et un
anneau nœthérien avec un idéal $\mathfrak{m}$ : la clause finale « si $A$
local » suppose la seconde lecture, mais « local » n'est pas biffé dans la
première ligne.}

Prenons $\mathfrak{m}$ premier, on \uncertain{conclut} :

\textbf{Corollaire 2 :} Soit $\mathfrak{p}$ un idéal premier de $A$ local
nœth., $r = \dim A_{\mathfrak{p}}$, alors il existe $x_1, \dots, x_r \in
\mathfrak{p}$ tels que les $x_i$ \struck{fassent part} définissent un syst. de
paramètres de $A_{\mathfrak{p}}$, et forment partie d'un système de
paramètres de \uncertain{$A / \mathfrak{p}$}.
\note{le dernier symbole se lit $A / \mathfrak{p}$ sur le feuillet ; on
attendrait $A$.}

\marginal{Corollaire. Conditions équivalentes sur $x_1, \dots, x_p \in
\mathfrak{m}_A$, $A$ local nœth. : (i) $\operatorname{ht}(\sum A x_i) = p$
(\uncertain{comme} \uncertain{toujours} $\leq$) ; (ii) \struck{Pour tout}
$x_1$ ne s'annule sur aucune comp. irréd. de $\operatorname{Spec}(A)$, …,
$x_n$ ne s'annule sur aucune comp. irréd. de $\operatorname{Spec} A / \sum_{i
< n} x_i A$.}
\note{écrit de bas en haut dans la marge de gauche, sur toute la hauteur du
feuillet ; la fin de (ii), en très petits caractères, est lue d'après le
sens : le quotient et son indice de sommation sont incertains.}

\page{84}
\note{le feuillet entier est barré de traits diagonaux ; il se lit sous la
rature.}

\textbf{$\partial$-Catégorie} $\mathcal{H}$ : catégorie additive graduée, i.e.
où les $\operatorname{Hom}(X, Y)$ sont des groupes gradués (la composition des
morphismes étant additive sur les degrés), satisfaisant aux axiomes qui vont
suivre. On désigne par $\mathcal{H}_{(0)}$ la catégorie $\mathcal{H}$ où les
hom sont réduits au degré $0$.

\textbf{Axiome $\partial 1$.} Pour tout $X \in \operatorname{Ob} \mathcal{H}$,
il existe un \struck{$X(1)$} \struck{$T(X) = X(1)$} $\in \operatorname{Ob}
\mathcal{H}$ \add{noté $X(1)$ ?} et un isomorphisme fonctoriel (sur
$\mathcal{H}_{(0)}$) en $X$
\[
  \operatorname{Hom}_1(X, Y) \xrightarrow{\ \sim\ } \operatorname{Hom}_0(X, Y(1)) .
\]
\note{le membre de droite est surchargé : le $(1)$ de $Y(1)$ est écrit
par-dessus une première version.}

Donc il y a un objet \ill{} $i_Y \in \operatorname{Hom}_1(\struck{Y(1)}, Y)$
ayant une propriété universelle. Évidemment $(Y(1), i_Y)$ \uncertain{est}
\uncertain{déterminé} \ill{} à \struck{Hom-}isomorphisme près.
\note{ces deux phrases sont encadrées et reliées par un trait à la ligne de
l'axiome ; l'argument de $\operatorname{Hom}_1$ est repassé et ne se décide
pas.}

\add{il existe} un objet $X(1)$ de $\mathcal{H}$ et un isomorphisme de degré
$1$
\[
  \alpha_X : X(1) \xrightarrow{\ 1\ } X
\]
et un objet $X(-1)$ de $\mathcal{H}$ et un isomorphisme de degré $-1$
\[
  \beta_X : X(-1) \xrightarrow{\ -1\ } X .
\]
On en conclut, pour tout $Y$, des isomorphismes
\[
  \text{I}
  \begin{cases}
    \operatorname{Hom}_n(X, Y) \xrightarrow{\ \sim\ } \operatorname{Hom}_{n+1}(X(1), Y) \\
    \operatorname{Hom}_n(Y, X(1)) \xrightarrow{\ \sim\ } \operatorname{Hom}_{n+1}(Y, X)
  \end{cases}
  \qquad
  \text{II}
  \begin{cases}
    \operatorname{Hom}_n(X, Y) \xrightarrow{\ \sim\ } \operatorname{Hom}_{n-1}(X(-1), Y) \\
    \operatorname{Hom}_n(Y, X(-1)) \xrightarrow{\ \sim\ } \operatorname{Hom}_{n-1}(Y, X) \struck{\ill{}}
  \end{cases}
\]
\marginal{$= \operatorname{Hom}_{n-1}(X, Y(1))$}
\note{en marge de droite, à hauteur de la première ligne de I, surchargé.}

\page{86}
\note{ce feuillet est le verso auquel renvoie la marge du feuillet 87 (« cf
verso ») : il porte le lemme dont la Remarque commencée au feuillet 88 se
sert. Ordre de lecture : 88, 87, puis 86.}

\struck{\ill{}} \textbf{Lemme.} Soient $A$ un anneau, $B$ une algèbre finie
sur $A$, $M$ un $B$-module de type fini.

(i) si $M_{[A]}$ est $(S_k)$, il en est de même de $M$.

(ii) Supposons que pour \struck{\ill{}} idéaux premiers $\mathfrak{q},
\mathfrak{q}'$ de $B$, au-dessus \add{d'un même} $\mathfrak{p}$ de $A$, on ait
$\dim M_{\mathfrak{q}} = \dim M_{\mathfrak{q}'}$. Alors si $M$ est $(S_k)$ sur
$B$, il est $(S_k)$ sur $A$.
\note{au-dessus de « au-dessus d'un même $\mathfrak{p}$ », en interligne :
« $\mathfrak{q} = \ill{}$ ».}

On peut supposer $A \subset B$, \struck{$\operatorname{Spec} B =$}
$\operatorname{Supp} M = \operatorname{Spec} B$.

Dém. (i) Soit $\struck{\mathfrak{q}}\ \mathfrak{p} \in \operatorname{Spec}
\struck{B}\ A$, \struck{si $\dim M_{\mathfrak{p}} \leq k$ \ill{} \ill{}}
\struck{l'hypothèse implique que $\dim_{A_{\mathfrak{q}}} M_{\mathfrak{q}} =
\dim_{B_{\mathfrak{p}}} M_{\mathfrak{p}}$, or \ill{} \ill{}} \struck{on a
$\operatorname{prof}$ [\ill{} $\dim_{A_{\mathfrak{q}}} M_{\mathfrak{q}} \geq
\operatorname{Sup}_{\mathfrak{p} \to \mathfrak{q}} \dim_{B_{\mathfrak{p}}}
M_{\mathfrak{p}}$] et on} on a :
\[
  \operatorname{prof}_{A_{\mathfrak{q}}} M_{\mathfrak{q}}
  = \operatorname{Inf}_{\mathfrak{p} \to \mathfrak{q}}
    \operatorname{prof}_{B_{\mathfrak{p}}} M_{\mathfrak{p}} ,
  \qquad
  \dim_{A_{\mathfrak{q}}} M_{\mathfrak{q}}
  = \operatorname{Sup}_{\mathfrak{p} \to \mathfrak{q}}
    \dim_{B_{\mathfrak{p}}} M_{\mathfrak{p}} ,
\]
donc si $\mathfrak{p} \to \mathfrak{q}$,
\[
  \operatorname{prof}_{B_{\mathfrak{p}}} M_{\mathfrak{p}}
  \geq \operatorname{prof}_{A_{\mathfrak{q}}} M_{\mathfrak{q}}
  \geq \operatorname{Inf}(k, \dim_{A_{\mathfrak{q}}} M_{\mathfrak{q}})
  \geq \operatorname{Inf}(k, \dim_{B_{\mathfrak{p}}} M_{\mathfrak{p}}) ,
  \quad \text{cqfd.}
\]
\note{dans l'énoncé, $\mathfrak{q}$ est un premier de $B$ au-dessus de
$\mathfrak{p}$ dans $A$ ; dans la démonstration, à partir de la ligne
biffée « l'hypothèse implique… », c'est $\mathfrak{p}$ qui est dans $B$ et
$\mathfrak{q}$ dans $A$, la flèche $\mathfrak{p} \to \mathfrak{q}$ notant
« au-dessus de ». L'échange est sur le feuillet et n'a pas été régularisé.}

(ii) On a ici, grâce à l'hyp. \uncertain{sur} $M_{\mathfrak{p}}$ \ill{} :
$\dim_{A_{\mathfrak{q}}} M_{\mathfrak{q}} = \dim_{B_{\mathfrak{p}}}
M_{\mathfrak{p}}$. Or pour $\mathfrak{p}$ convenable sur $\mathfrak{q}$, on
aura \ill{} \struck{\ill{}} $\operatorname{prof}_{A_{\mathfrak{q}}}
M_{\mathfrak{q}} = \operatorname{prof}_{B_{\mathfrak{p}}} M_{\mathfrak{p}}$
(i.e. \ill{} $\geq$) \ill{}.

\struck{Pour \ill{} \ill{} \ill{} \ill{} $M_{\mathfrak{q}}$ \ill{} \ill{}
\ill{} on a \ill{} \ill{} \ill{}}
\note{le bas du feuillet, une dizaine de lignes, est hachuré au point de ne
plus se lire ; seuls surnagent $M_{\mathfrak{q}}$, « on a », un petit
schéma $B$ sur $A$ avec $M$ à côté, et en bas à gauche
« $\operatorname{prof} \geq \operatorname{Inf}(k, \dim)$ ». En marge de
gauche, biffé : $\mathfrak{p} \subset \mathfrak{q}$.}

\page{87}
\note{suite directe du feuillet 88 : la phrase commencée au bas de 88 se
termine ici.}

il s'ensuit que $\mathcal{C}$ est \uncertain{cohérent}. D'ailleurs,
$\mathcal{C}$ satisfait $(S_2)$ en tout point de $Z$, \struck{en \ill{}
\ill{}} \struck{\uncertain{Comme} \ill{} $\mathcal{M} \subset X$, \ill{}
$\Gamma(V, \mathcal{C}) \to \Gamma(V \cap U, \mathcal{C})$ \uncertain{est}
\ill{}} \struck{\ill{}} par construction $\mathcal{C} \xrightarrow{\ \sim\ }
i_* i^* \mathcal{C}$. Comme $\mathcal{C} \simeq \mathcal{B}$ \ill{}
\uncertain{sur} $U$, \struck{\ill{}} et que $\mathcal{B}_x$ est une algèbre
finie et \uncertain{normale} sur $\mathcal{O}_{X, x}$, donc $(S_2)$ comme
\uncertain{anneau}, et par suite comme \uncertain{module}, il s'ensuit que
$\mathcal{C}$ \uncertain{est} $(S_2)$ \struck{comme Module}
\uncertain{partout}. \struck{\ill{}} \ill{} \uncertain{que}
\uncertain{l'anneau} $\mathcal{C}$ \ill{} \ill{} $\mathcal{C}$ est $S_2$. Bien
entendu $A \subset \mathcal{C} \subset K$. Reste : prouver que $\mathcal{C}$
est \struck{normal} $(R_1)$. Soit $\mathfrak{p}$ un idéal premier de
$\mathcal{C}$ tel que $\dim \mathcal{C}_{\mathfrak{p}} = 1$, prouvons
$\mathcal{C}_{\mathfrak{p}}$ régulier. Soit $\mathfrak{q}$ l'idéal induit par
$\mathfrak{p}$ sur $A$, alors [$A$ \uncertain{est} \uncertain{quotient} d'un
régulier] on a $\dim A_{\mathfrak{q}} = 1$, donc ou bien $A_{\mathfrak{q}}$
est normal, donc $A_{\mathfrak{q}} = B_{\mathfrak{q}} = B_{\mathfrak{p}}$
\struck{\ill{}} normal, ou bien $A_{\mathfrak{q}}$ non normal, i.e. on est en
$\mathfrak{p}_i$, mais alors \ill{} \ill{} \uncertain{par} construction
$B_{\mathfrak{p}}$ normal, \struck{\ill{}} $\mathcal{B}_{\mathfrak{p}}$ est
normal comme localisé de $B_{\mathfrak{q}}$. Cela achève la démonstration.

\marginal{[\uncertain{A} \ill{} \ill{} régulier \ill{} $\dim A_{\mathfrak{q}}
= \dim B_{\mathfrak{p}}$ \uncertain{et} $\operatorname{prof} \mathcal{B}$
\uncertain{induit} \ill{}] cf verso}
\note{écrit en oblique dans la marge de gauche, à hauteur de « Reste :
prouver » ; « cf verso » renvoie au lemme du feuillet 86.}

\struck{$A$ anneau (\uncertain{local nœth.} \ill{}), $B$ \uncertain{alg.}
finie sur $A$, $M$ module de type fini sur $B$, \ill{} \ill{}. Alors $M$ est
$(S_k)$ sur $A$ ssi il est $(S_k)$ sur $B$. \struck{Soit} \ill{} \ill{}
\uncertain{Supposons} $A \subset B$, $\operatorname{Supp} M =
\operatorname{Spec} B$, \ill{} $f(x) = y$. Soit $y \in Y$, …
$\operatorname{prof} M_y = \operatorname{Inf}_{x \mapsto y} \operatorname{prof}
M_x$ \ill{} Si $M$ sur $B$ est $(S_k)$, \ill{} \ill{} $\operatorname{prof} M_y
\geq \operatorname{Inf}$ \ill{}}
\note{le bas du feuillet : première rédaction, hachurée, du lemme du feuillet
86, avec $f : X \to Y$ écrit verticalement à gauche et, en insertion
au-dessus de la dernière ligne, « il \ill{} \ill{} \uncertain{dimension} ».
Cette version énonçait une équivalence là où le lemme retenu ne donne que
deux implications sous hypothèse.}

\page{88}

\textbf{Remarque.} Soit $A$ un anneau intègre quotient d'un anneau régulier,
$K$ son corps des fractions, \struck{$L$ une extension finie de $K$}. Pour que
la clôture intégrale de $A$ soit de type fini, il f. et s. que $A$ satisfasse
les conditions suivantes :

a) l'ensemble des $\mathfrak{p}_i \in \operatorname{Spec} A$ tel que
$A_{\mathfrak{p}_i}$ soit de dim $1$ non régulier est fini

b) Pour tout tel $\mathfrak{p}_i \in \operatorname{Spec}(A)$, \struck{tel que
$\dim A_{\mathfrak{p}} = 1$,} la normalisée de $A_{\mathfrak{p}}$ est finie
\struck{sur} $A_{\mathfrak{p}}$.

Dém. Nécessité évidente [a) car le lieu \struck{régulier} non normal de
$\operatorname{Spec} A$ \uncertain{est} fermé, et les $\mathfrak{p}_i$
correspondent aux composantes irréd. de dim $1$ \uncertain{de celui-ci}].

Suffisance. \struck{Soit $B \subset K$ une \uncertain{alg.} finie sur $A$, tel
que \struck{Les} $B_{\mathfrak{p}_i}$ soient normaux [existe grâce à a),
b)].} Comme $A$ est quotient d'un régulier, l'ens. des $x \in
\operatorname{Spec} A = X$ tel que $A_x$ soit $S_2$ est ouvert, et contient
\struck{les} \ill{} un ouvert \ill{} \ill{}, soit $X_f$.
\note{en interligne au-dessus de « contient … soit $X_f$ » : « (utilisant
a), et le critère de Serre, i.e. \uncertain{que} \uncertain{les}
$\mathfrak{p}_i$ \uncertain{sont} \ill{}) », et au-dessus de $X_f$ :
\uncertain{normal}.}
Il existe \struck{\ill{}} un sous-anneau $B \supset A$ de $K$, fini sur $A$,
\struck{\ill{} $B_f$ \ill{}} tel que $B_{\mathfrak{p}_i}$ soit normal [utilise
a) et b)]. Soit $Z = \operatorname{Supp} B / A$, \struck{$Z$} donc $Z \subset
X - X_f$, et $Z$ est de codim $> 1$. Soit $U = X - Z$, $i : U \to X$
l'immersion canonique, $\mathcal{B} = \widetilde{B}$ faisceau d'alg. sur $X$,
considérons \struck{$\mathcal{B}$} $\mathcal{C} = i_* i^*(\mathcal{B})$. Comme
le \ill{} \uncertain{anneau} : $\mathcal{B}$ \uncertain{est} \ill{} $X$, et
que $Z$ est de codim $\geq 2$ dans $X$, \ill{} \struck{$A$}
\uncertain{partout} \uncertain{dans} \ill{},
\note{la phrase se poursuit au feuillet 87 : « il s'ensuit que $\mathcal{C}$
est cohérent ». Au-dessus de « tel que $B_{\mathfrak{p}_i}$ » : une coche et
« $\operatorname{ht}$ », d'une autre encre.}

\page{89}
\note{brouillon pour le lemme du feuillet 86, à l'encre, barré de traits
diagonaux. Sous l'encre, une couche antérieure au crayon, en anglais, sur le
groupe de Picard et la profondeur ; elle est transcrite à la suite.}

$A \subset B$, $M$ \struck{\ill{}}. $\operatorname{prof} M_y \struck{\leq}
\struck{\operatorname{Inf}} \operatorname{prof} M_z \geq
\operatorname{Inf}_{z \mapsto y} \struck{\dim(M_z, k)} \operatorname{Inf}(\dim
M_z, k)$.

\struck{\ill{}} $B$ \uncertain{normal}, $M$ \struck{$\leq$ \ill{}
$\operatorname{Inf}$}. $S_2$ comme \uncertain{anneau}, mais pas comme
module.
\note{« mais pas comme module » est souligné d'un trait ondulé.}

$B \subset A^n$, $B / \mathfrak{f} B \subset (A / \mathfrak{f} A)^n$ ;
$B / \mathfrak{f} B \to (A / \mathfrak{f} A)^n$ ; $B / \mathfrak{f} B \simeq
\mathcal{O}_{X_0}$ \ill{}. $B$ \uncertain{est} \uncertain{C.M.}
\uncertain{aussi} \ill{}. \struck{\ill{}}

$M$ $S_k$ $\overset{?}{\Longrightarrow}$ $M_{[A]}$ est $(S_k)$.

$\mathfrak{p} \in A$, $\dim A_{\mathfrak{p}} \leq k \Rightarrow \dim
B_{\mathfrak{p}} \leq k$ car $M_{\mathfrak{p}}$ est C.M. sur
$B_{\mathfrak{p}}$, donc $M_{\mathfrak{p}}$ est C.M. \struck{\ill{} sur
\ill{}} \struck{\ill{} sur $A$}.
\note{« $\mathfrak{p} \in A$ » est encadré.}

$M_{\mathfrak{q}}$, $\dim M_{\mathfrak{q}}$, $\operatorname{prof}
M_{\mathfrak{q}}$ ; $B$ est C.M. \ill{} \ill{} ; $B_{\mathfrak{q}}$,
$A_{\mathfrak{p}}$ ; $\dim B_{\mathfrak{q}} = \operatorname{prof}$ \ill{}
\[
  \operatorname{prof} M_{\mathfrak{p}} \geq \operatorname{Inf}(k, \dim M_{\mathfrak{p}})
  \iff
  \operatorname{prof} \hat{M}_{\mathfrak{p}} \geq \operatorname{Inf}(k, \dim \hat{M}_{\mathfrak{p}}) ,
  \qquad
  \operatorname{prof} \hat{M}_{\mathfrak{q}_i} \geq \operatorname{Inf}(k, \dim \hat{M}_{\mathfrak{q}_i}) .
\]

\note{la couche au crayon, en anglais, écrite dans l'autre sens du feuillet
(en colonne, à droite et à gauche) :}

$\varphi : \operatorname{Pic}(X') \to \operatorname{Pic}(Y')$ ;
$\pi_1(X')$, $\pi_1(Y')$, $H^i(\ldots)$.

1) if $\forall x \in X' - Y'$, $\dim \bar{x} = 1 \Rightarrow \operatorname{depth}
A_x \geq 3$, $\varphi$ is injective.

2) if $\forall x \in X' - Y'$, $\dim \bar{x} = 1 \Rightarrow
\operatorname{depth} A_x \geq 4$ and if $\operatorname{Pic}(\operatorname{Spec}
A_x - \{x\}) = 0$ $\Rightarrow$ $\varphi$ is bij.

complete intersection, $\dim \geq 4$ ; $\operatorname{Pic}(X - a) = 0$ ;
$\mathfrak{a} = \mathfrak{f} A + \mathfrak{g} A$ ; $\mathfrak{a}(\mathfrak{m})$.
\note{ces derniers fragments sont dans la marge de gauche, au crayon, en
partie sous l'encre ; leur liaison avec 1) et 2) ne se lit pas.}

\section{IV 12. Fibres des morphismes plats de présentation finie}
\note{les mots sont de sa main, en haut à droite de la chemise orange qui
forme le feuillet 90 : « IV \uncertain{12} » à l'encre, le titre au crayon.
Le second chiffre est repassé et pourrait aussi se lire 82. La chemise ne
porte rien d'autre.}

\page{91}
\note{brouillon rapide, encre par-dessus crayon, avec un croquis à gauche :
deux disques $X'$ et $X''$ marqués $Z'_0$, $Z'$, $Z''_0$, $Z''$ autour d'un
centre, reliés par une courbe $D$ ; les lignes du texte sont lues dans
l'ordre où elles se suivent, les insertions ramenées à leur place.}

\textbf{Un exemple} [sur une \uncertain{conj.} \uncertain{log.} des \ill{}
\ill{} \ill{}]
\marginal{Où le \struck{\ill{}} des \uncertain{ouverts} des pts non Coh. Mac.
fait un saut dans la mauvaise direction…}
\note{en accolade, en haut à droite, d'une encre plus noire que le titre.}

\struck{$A$} $X'$ schéma local \add{intègre} \struck{de prof \ill{} \ill{}
\ill{}} de dim $n'$ \uncertain{(centre $x$)}, \struck{contenant une partie
\ill{} irréductible $Z_1 \ni x$,} telle que aux pts de $Z_1$, $X'$ ne soit
\struck{\ill{}} \uncertain{Cohen-Macaulay}, \ill{} \uncertain{donc} $Z'_0$
l'un des pts \ill{} \uncertain{il} \uncertain{n'est} \uncertain{pas} Coh. M.
\struck{$\operatorname{codim}_x Z_1 =$ \ill{}}
\note{au crayon, en marge de droite : « i.e. \uncertain{coprof} 1 » ; en
marge de droite encore, entouré : $Z'_0 \neq \{x\}$, et sous lui une ligne
noircie.}

$X''$ schéma local \add{intègre} de dim \uncertain{$n''$}, $Z''_0$ l'un des pts
\ill{} il n'est pas Coh. Mac.
\note{entouré, en interligne : « $n'' > n'$ ».}

$D$ un \struck{\ill{}} \uncertain{schéma} \uncertain{local} \uncertain{intègre}
\struck{fermé dans $X$} de dim $1$ \uncertain{et} \ill{}, $D'$ \struck{\ill{}}
(immersions fermées)
\[
  u' : D \to X' , \qquad u'' : D \to X'' , \qquad u'(D) \not\supset Z'_0 ,
\]
\[
  X = X' \coprod_D X'' .
\]
\note{« $u'(D) \not\supset Z'_0$ » est entouré, et souligné deux fois.}

$\sim$ \uncertain{identifie} $X'$ et $X''$ ; \uncertain{donc} \ill{} \ill{}
\uncertain{schémas} \uncertain{finis}, \uncertain{donc} $Z'_0, Z''_0 \to x$.
$D$ : un \struck{\ill{}} \uncertain{sous-schéma} \uncertain{fermé}
\struck{\ill{}} de $X$.

$f_0 \in \struck{\ill{}} (\mathcal{O}_{D, x})$, $f_0 \neq 0$ ; les
\uncertain{relèvements} $f'$ \uncertain{et} $f''$, d'où $f \in \mathcal{O}_{X,
x}$ \ill{} \uncertain{qui} \uncertain{est} \uncertain{nul} \ill{} \ill{}
\[
  V(f) \simeq V(f') \coprod_{V(f_0)} V(f'') .
\]
\note{en bas à gauche, l'algèbre du recollement : $A' \times_B A''$ ;
$C \to A' \to B$, $C \to A'' \to B$ ; $f = (g', g'')$, $(g', g'') \in C$ ;
$g' \in \mathfrak{f}'$, $g'' \in \mathfrak{f}''$ ; $A' \times A'' \to B \to
0$. Les relations entre ces lignes ne sont pas écrites.}

\page{92}

[il \uncertain{pourra} \ill{} \ill{} \uncertain{fini} \ill{} \ill{} $x$]
\note{au crayon, en haut à droite.}

$V(f')$ et $V(f'')$ sont sans idéaux immergés.
\note{entouré. À gauche, isolé : « $\neq$ ? ».}

[les \uncertain{pts} = les \uncertain{pts} intègres]
\[
  (V(f') \cap Z'_0 =)\ V(f') \cap Z'_0 = \{x\} , \qquad
  (V(f'') \cap Z''_0 =)\ V(f'') \cap Z''_0 \subset \{x\} .
\]
\note{les deux lignes sont entourées ; les parenthèses de gauche sont sur le
feuillet, une première rédaction non biffée.}

Il s'ensuit que $V(f)$ est sans idéaux immergés, [et \ill{} \ill{} si \ill{}
\ill{}], et que $V(f)$ est Coh. Mac. en dehors de $\{x\}$. \struck{Donc
\ill{} pt} Donc l'un des pts où $X$ n'est pas Coh. Mac. \uncertain{est}
\[
  Z = \struck{\ill{}} Z'_0 \cup Z''_0 \cup D ,
\]
\uncertain{qui} \uncertain{donc} \ill{} ; $X - V(f)$ \uncertain{est} de codim
$\leq \operatorname{codim}_{X'} Z'_0 \leq n'$ \ill{} [car $Z'_0 \neq
\emptyset$, $Z'_0 \neq \{x\}$], \struck{et qui donc} \uncertain{alors}
\uncertain{que} le codim de $V(f)$ \uncertain{en} $V(f) \cap Z$ [i.e. l'un des
pts où $V(f)$ n'est \struck{\ill{}} \uncertain{il} \uncertain{n'est} pas Coh.
Mac.] \uncertain{est} de codim $= \struck{\ill{}} \dim V(f) = n'' - 1 \geq
n'$.
\note{la borne « $\leq n'$ » est suivie d'un signe surchargé qui pourrait être
« $-1$ » ; l'inégalité finale $n'' - 1 \geq n'$ est celle qui fait
« sauter dans la mauvaise direction » la codimension du lieu non
Cohen-Macaulay quand on passe de $X$ à $V(f)$.}

\page{93}

Pour \uncertain{faire} la construction \add{d'un exemple} : \uncertain{un}
\uncertain{pt} d'un schéma \add{intégralement} \struck{\ill{}} projectif
\uncertain{normal} $S'$ \add{dans un espace projectif \ill{}}, \struck{\ill{}
\ill{}} qui n'est pas Cohen-Macaulay [\uncertain{que} \ill{} de dim $d \geq
3$] \struck{\ill{} [\ill{} \ill{}]}, et pour $X'$ le schéma localisé en le
sommet du cône projetant de \struck{$X'$, $X'$} \add{$S'$}, \uncertain{qu'on}
\uncertain{peut} \uncertain{prendre} projectivement normale. Donc $n' = d + 1
\geq 4$. \struck{Le \ill{} $X''$ \ill{}} \uncertain{soit} \uncertain{que}
\struck{\ill{}} $Z'_0 \neq \emptyset$, $Z'_0 \neq \{x\}$ est
\uncertain{vérifié}. On prend $X'' = \struck{\ill{}}$ localisé en l'origine de
$V(k^{n''})$, \uncertain{i.e.} $n'' > n' = d + 1$. Il existe \uncertain{alors}
un \uncertain{morphisme} \add{d'\uncertain{immersion} \uncertain{fermée}
$u'$} de \uncertain{$D = \operatorname{Spec} k[t]_{(t)}$} dans $X'$, qui soit
tel que $u'(D) \not\supset Z'_0$ [prendre une \uncertain{générisation}
\uncertain{non} \uncertain{générique}], et on prend une immersion $u'' : D \to
X''$ (\uncertain{par} une \uncertain{droite} \uncertain{à} \uncertain{l'origine}
\uncertain{fixée}). On choisit $f_0$ \uncertain{quelconque}, puis $f'$ et
$f''$ de façon à vérifier $V(f') \cap Z'_0 = \{x\}$, \uncertain{et}
d'ailleurs $V(f'') \cap Z''_0 \subset \{x\}$ est \uncertain{trivial}. Le
\uncertain{caractère} \uncertain{de} $X'$ et $X''$ implique que $V(f')$ et
$V(f'')$ n'ont pas d'idéaux \uncertain{immergés} \ill{}.
\note{le mot « projectivement normale » (au féminin) et « le sommet du cône
projetant » sont insérés en interligne ; l'écriture de $D$ porte un indice
sous $k[t]$ qui se lit mal, $(t)$ d'après le sens. Le feuillet s'arrête sur
cette phrase ; le verso, feuillet 94, est blanc.}

\page{95}

\marginal{Compl\ill{} \uncertain{Proposition} \ill{} \ill{} des
« \uncertain{conditions} » « \uncertain{signalétiques} ». Modules
[\uncertain{généralisation} : un $B$-module $M$]}
\note{écrit en oblique dans la marge de gauche, sur deux ou trois lignes
soulignées, le mot « Modules » plus gros que le reste ; la nature du
complément annoncé ne se lit pas.}

\textbf{Proposition.} Soient $K$ un \uncertain{anneau}, $A, B, C$ des
\uncertain{algèbres} \uncertain{graduées} sur $K$, à degrés positifs, et
\uncertain{réduites} à $K$ en degré $0$,
\[
  A \xrightarrow{\ \psi\ } B \xrightarrow{\ \varphi\ } C
\]
des homomorphismes. On suppose

a) Pour tout $i \geq 0$, $\varphi^i : B^i \to C^i$ est surjectif et admet un
inverse à droite \uncertain{$K$-linéaire}, soit $u^i : C^i \to B^i$.

(1) Alors les conditions b) et b') suivantes sont équivalentes :

b) \struck{\ill{}} $\operatorname{Ker} \varphi = \psi(A^+) B$

b') Pour tout $n \geq 0$, l'homomorphisme
\[
  v^n : \sum_{i + j = n} A^i \otimes_K C^j \longrightarrow B^n
\]
\struck{\ill{}} \uncertain{défini} \uncertain{par} $v^n(a^i \otimes c^j) =
\psi(a^i) \struck{\ill{}} u^j(c^j)$ est surjectif.

[Donc (b') ne dépend pas du choix des $u^j$].

\marginal{N.B. [sans a)] : $A$ \uncertain{augmenté}, $A'$, \ill{} ; b)
\uncertain{équivaut} : $\operatorname{Ker} \varphi = \psi(A') B$ ; b1)
$\operatorname{Ker} \varphi^1 = \psi(A^1)$ ; b2) $\operatorname{Ker}
\varphi$ \ill{} $\operatorname{Ker} \varphi^n$ ; \ill{}}
\note{écrit en oblique dans la marge de gauche, à hauteur de b) ; « b2) »
est repassé et une flèche relie ses deux membres.}

\page{96}

(2) Supposons \uncertain{que} \ill{} \ill{} \ill{}

b'') $v^n$ est bijectif.

Alors $\psi$ est injectif. \struck{Alors} De plus, les $B^i$ sont plats
[\uncertain{resp.} \struck{\ill{}} projectifs] sur $K$ \uncertain{ssi} les
$A^i, C^i$ le sont.

\struck{Réciproquement, supposons que les $A^i, B^i, C^i$ projectifs de type
fini (= plats de présentation finie), pour tout}
\note{biffé d'un trait, puis repris tel quel dans (3) ; une accolade en marge
rattache le bloc biffé à ce qui suit. La parenthèse « = plats de présentation
finie » est de lecture incertaine.}

(3) $x \in \operatorname{Spec} K$, \uncertain{soient}
\[
  P_{A, x}(t) = \sum_i \operatorname{rg}_{k(x)}(A^i(x))\, t^i \in \mathbb{Z}[[t]]
\]
et \uncertain{définissons} de même $P_{B, x}(t)$, $P_{C, x}(t)$. Alors
\uncertain{on a} pour tout $x$ [\uncertain{en} \uncertain{grandeur} \ill{}]
\[
  (\ast) \qquad P_{B, x} \ll P_{A, x} P_{C, x}
\]
[\uncertain{inégalité} \uncertain{coefficient} par \uncertain{coefficient}].
Si (b'') est vérifié, alors on a $=$.

c) $P_{B, x} = P_{A, x} P_{C, x}$.
\note{c) est écrit en marge de gauche, au bout d'un trait qui le relie à la
phrase suivante.}

Si les $A^i, B^i, C^i$ sont plats de prés. finie [= projectifs de type fini],
alors c) implique aussi b''), \struck{\ill{}} \struck{Donc}

\marginal{N.B. \uncertain{Ce} \uncertain{dernier} \uncertain{résultat}
\uncertain{que} (b'') \uncertain{ne} \uncertain{dépend} \ill{} \uncertain{Alors}
\uncertain{les} \uncertain{des} deux \struck{\ill{}} \uncertain{parties}
\ill{} [si les $A^i, B^i$ \uncertain{sont} de type fini, \ill{} \ill{}
\uncertain{seulement} \uncertain{de} \ill{} \uncertain{en} \ill{} d'un \ill{}
\uncertain{de} \ill{}] (3) (\uncertain{continue} \ill{}, 3°) grâce à
Nakayama] +}
\note{écrit en oblique dans la marge de gauche sur toute la hauteur du
feuillet ; le « (3) » entouré renvoie au point (3). Le nom de Nakayama est sûr, le reste
de la note ne se lit qu'en partie.}

\page{97}

\textbf{Corollaire 1.} Conditions équivalentes [\uncertain{moyennant} a)] :
\note{« moyennant a) » est en interligne au-dessus, rattaché par un trait.}

(i) $v^*$ bijectif, les $B^i$ projectifs de type fini

(ii) $v^*$ bijectif, les $A^i, C^i$ projectifs de type fini

(iii) les $A^i, B^i, C^i$ projectifs de type fini, \struck{et}
$P_{B, x} = P_{A, x} P_{C, x}$ pour tout $x$

[N.B. si $K$ est local, il suffit de regarder en le pt fermé].

\textbf{N.B.} Nous \uncertain{appellerons} \ill{} l'ensemble des conditions
a), b'') \uncertain{comme} « bonnes », si de plus les $B^i$ sont projectifs
de type fini, on dira « très bonnes ».

\textbf{Question.} Sur les « bonnes conditions » \uncertain{peut-on}
\uncertain{prouver} :

(i) $A$ plat$/K$ $\Longrightarrow$ $C$ plat$/B$

(ii) $C$ plat$/K$ $\Longrightarrow$ $B$ plat$/A$

(iii) $A$ et $C$ plats$/K$ $\Longrightarrow$ $B$ plat$/K$
\note{à gauche de (iii), un mot isolé, « dans » ou « donc », qui ne se
rattache à rien.}

\page{98}

\textbf{Corollaire 2.} Supposons que $C \simeq K[t_1, \dots, t_r]$
[\uncertain{avec} \uncertain{des} \uncertain{degrés} $> 0$], les $t_i$
homogènes. Alors (a, b'') équivalent à ceci :
\[
  \begin{cases}
    B \simeq A[t_1, \dots, t_r] \\
    \psi \simeq \text{hom. can. } A \to A[t_1, \dots, t_r] \\
    \varphi \simeq \text{hom. \uncertain{projection} \uncertain{canonique} }
      A[t_1, \dots, t_r] \to \underbrace{A / A^+}_{K}[t_1, \dots, t_r]
  \end{cases}
\]
[Le choix d'un isom. $B \simeq A[t_1, \dots, t_r]$ équivaut : \uncertain{celui}
d'un hom. \uncertain{d'alg.} $C \xrightarrow{\ u\ } B$ tel que $\varphi u =
\operatorname{id}_C$, \uncertain{ou} \uncertain{encore} d'un hom.
\uncertain{$K$-linéaire} \uncertain{qui} \uncertain{relève} $B^+ \to C^+ /
C^{+2}$].
\note{« $C \xrightarrow{u} B$ » est souligné ; devant lui un mot biffé.
Le feuillet, et la cote, s'arrêtent ici.}

\end{document}
